\documentclass[11pt,a4paper]{article}
\usepackage{isabelle,isabellesym}
\usepackage{amsmath,amssymb}
\usepackage{url}
\usepackage{pdfsetup}

\urlstyle{rm}
\isabellestyle{it}

\begin{document}

\title{Steiner's Line Theorem in Isabelle/HOL}
\author{Arthur Freitas Ramos \and David Barros Hulak \and Ruy Jose Guerra Barretto de Queiroz}

\maketitle

\begin{abstract}
We formalize Steiner's line theorem in Isabelle/HOL.  For a nondegenerate
triangle and a point on its circumcircle, the reflections of that point in the
three sidelines are collinear, and the resulting line passes through the
orthocenter.  The development uses complex coordinates and imports the
published Wallace--Simson line entry for its perpendicular-foot and
collinearity infrastructure.
\end{abstract}

\section{Introduction}

Steiner's line theorem is the reflection companion to the Wallace--Simson
line theorem~\cite{CoxeterGreitzer1967}.  Its classical setting goes back to
Steiner's work on triangle and circle geometry~\cite{Steiner1826}; a modern
annotated account of related Steiner theorems is given by Ehrmann
~\cite{Ehrmann2004}.  Let $ABC$ be a nondegenerate
triangle and let $M$ lie on its circumcircle.  Reflect $M$ in the sidelines
$BC$, $CA$, and $AB$.  The
three reflected points lie on one line, called the Steiner line of (M) with
respect to (ABC), and this line passes through the orthocenter of (ABC).

The formalization imports the AFP entry \emph{The Wallace--Simson Line
Theorem in Isabelle/HOL}~\cite{Simson-AFP}.  Each reflection is defined as
twice the corresponding perpendicular foot minus (M), so the collinearity
part is obtained by applying a homothety to the Simson line.  The
orthocenter incidence is then proved by a direct complex-coordinate identity.

\section{Main results}

The theorem \isa{steiner\_line} proves collinearity of the three reflected
points.  The stronger theorem \isa{steiner\_line\_through\_orthocenter}
proves all three cyclic orthocenter/reflection collinearities.  This
redundant-looking statement is intentional: if $M$ is a vertex, one pair of
reflections can coincide, while another cyclic pair still determines the
Steiner line.  The lemma \isa{orthocenter\_is\_orthocenter} verifies the
geometric meaning of the coordinate definition: the three
vertex-to-orthocenter displacements are perpendicular to the opposite
sidelines.  Its algebraic statement only assumes the three common-circle
conditions; the nondegenerate-triangle and positive-radius hypotheses are
introduced by the main theorems.  The definitions
\isa{steiner\_reflect} and \isa{orthocenter} expose the concrete
complex-coordinate construction used by both results.

\input{session}

\bibliographystyle{abbrv}
\bibliography{root}

\end{document}
